\documentclass[11pt]{article}

\usepackage[a4paper,margin=1in]{geometry}
\usepackage{amsmath,amssymb}
\usepackage{slashed}
\usepackage{hyperref}

\title{Momentum-Space QCD Sum-Rule Setup for the $B_c(1P)$ Axial-Vector Mixing Angle}
\author{}
\date{}

\begin{document}

\maketitle

\section{Purpose and conventions}

We study the mixing between the two axial-vector $B_c(1P)$ basis states
associated with the ${}^3P_1$ and ${}^1P_1$ structures.  The current
implementation follows the Aliev et al. mixing-angle method, but the OPE is
evaluated in momentum space because both quark lines are heavy.

This note records the calculation implemented in
\texttt{BcMixingMomentum.wl}.  It is written so that every convention and
intermediate object can be checked independently before using the numerical
angle in a paper.

The quark content convention is
\begin{equation}
  B_c^+ = \bar b c .
\end{equation}

The two interpolating currents are taken as
\begin{align}
  J^A_\mu &= \bar b\,\gamma_\mu\gamma_5\,c,\\
  J^B_\mu &=
  i\,\bar b\,\sigma_{\mu\alpha}
  \frac{p^\alpha}{m_b+m_c}
  \gamma_5\,c,
\end{align}
where
\begin{equation}
  \sigma_{\mu\nu}=\frac{i}{2}[\gamma_\mu,\gamma_\nu].
\end{equation}
The explicit factor of $i$ in $J^B_\mu$ is retained.  This is important
because FeynCalc's \texttt{DiracSigma} already contains the conventional
$i/2$ commutator.

The factor $1/(m_b+m_c)$ is also essential.  Since
$[\bar Q Q]=3$ and $[p^\alpha]=1$, the unnormalized tensor current
$i\bar b\sigma_{\mu\alpha}p^\alpha\gamma_5c$ has mass dimension 4, while
$J^A_\mu$ has mass dimension 3.  The normalized current above makes both
currents dimension 3.

If the opposite flavor convention is used, for example $B_c^- = \bar c b$,
the sign of the quoted angle can change.  The absolute convention must
therefore be kept fixed when comparing to other papers.

\section{Mass-dimension check}

The quark field has dimension $3/2$, while $p^\alpha$, $m_b$, and $m_c$ have
dimension 1.  Therefore
\begin{equation}
  [J^A_\mu]=3,
  \qquad
  [J^B_\mu]=3
\end{equation}
only after the tensor current is normalized by $m_b+m_c$.

For currents with dimensions $d_i$ and $d_j$,
\begin{equation}
  [\Pi^{ij}_{\mu\nu}]
  =
  d_i+d_j-4,
\end{equation}
because of the $d^4x$ integration in the two-point function.  Thus the
normalized basis gives
\begin{equation}
  [\Pi^{AA}_1]=[\Pi^{AB}_1]=[\Pi^{BB}_1]=2.
\end{equation}
The spectral densities have the same mass dimension as the invariant
correlator,
\begin{equation}
  [\rho^{AA}]=[\rho^{AB}]=[\rho^{BB}]=2.
\end{equation}
In the present code the Borel moment is evaluated as
$\int ds\,e^{-s/M^2}\rho(s)$, so all three numerical moments carry one common
extra factor of $[ds]=2$.  This common convention does not affect the mixing
angle.

If the tensor current were left unnormalized, the dimensions would instead be
\begin{equation}
  [\rho^{AA}],\ [\rho^{AB}],\ [\rho^{BB}]
  =
  2,\ 3,\ 4,
\end{equation}
and the matrix entries in the angle formula would not be dimensionally
compatible.

\section{Implementation map}

The main code objects are:
\begin{itemize}
  \item \texttt{Correlator[channel, order]}:
  $\Pi_1^{ij}(p^2)$ after FeynCalc Dirac algebra and spin projection.
  \item \texttt{SpectralDensity[channel, ``pert'', s]}:
  $\rho_{\rm pert}^{ij}(s)$.
  \item \texttt{NumericBorelPi[channel, order, M2, s0]}:
  $\Pi^{ij}_{\rm order}(M^2,s_0)$.
  \item \texttt{NumericMixingAngleDegrees[M2, s0, ``total'']}:
  $\theta(M^2,s_0)$ from
  $\Pi_{\rm pert}^{ij}+\Pi_{G^2}^{ij}+\Pi_{G^3}^{ij}$.
\end{itemize}

The implemented OPE orders are
\begin{equation}
  \text{pert},\qquad G^2,\qquad \text{pertG2}=\text{pert}+G^2,
  \qquad G^3,\qquad \text{total}=\text{pert}+G^2+G^3.
\end{equation}
The implemented dimension-6 triple-gluon condensate
$\langle g_s^3G^3\rangle$ is the standard vacuum-averaged single-line
propagator term. Possible cross-line open-field $G^3$ terms need a separate
derivation.

\section{How to reproduce the calculation while reading this note}

This section is a practical roadmap for a new student.  The goal is not to
evaluate the final number immediately, but to check one layer at a time.  A
useful order is:
\begin{enumerate}
  \item Check the current definitions and mass dimensions.  In particular,
  verify that the factor \(p^\alpha/(m_b+m_c)\) in \(J^B_\mu\) makes the
  \(AA\), \(AB\), and \(BB\) correlators dimensionally compatible.
  \item Reproduce the Wick contractions for \(\Pi_{\mu\nu}^{AA}\),
  \(\Pi_{\mu\nu}^{AB}\), and \(\Pi_{\mu\nu}^{BB}\).  At this stage do not
  calculate any integrals; only check the order of gamma matrices and signs.
  \item Apply the spin--1 projector
  \[
    P^{\mu\nu}_{(1)}
    =
    \frac{1}{3}
    \left(g^{\mu\nu}-\frac{p^\mu p^\nu}{p^2}\right)
  \]
  and verify that a scalar invariant \(\Pi_1^{ij}\) is obtained.
  \item For the perturbative term, put the two internal quarks on shell and
  verify the three numerator functions \(N^{AA}\), \(N^{AB}\), and \(N^{BB}\)
  given below.  This is the most important algebraic check.
  \item Verify the phase-space factor
  \[
    \frac{N_c}{16\pi^2}
    \frac{\sqrt{\lambda(s,m_b,m_c)}}{s}
  \]
  and the physical threshold \(s_{\rm th}=(m_b+m_c)^2\).
  \item Perform the Borel integral
  \[
    \Pi^{ij}(M^2,s_0)
    =
    \int_{s_{\rm th}}^{s_0}ds\,e^{-s/M^2}\rho^{ij}(s)
  \]
  and only then insert the three Borel moments into the mixing-angle formula.
  \item Add condensates only after the perturbative result is stable.  For
  \(G^2\), keep the local single-line terms and the cross-line open-field
  term conceptually separate; otherwise it is easy to double count or miss a
  contribution.
\end{enumerate}

The Mathematica workflow mirrors these steps:
\begin{center}
\begin{tabular}{ll}
Hand calculation object & Mathematica check \\
\hline
current dimensions & \texttt{MixingMassDimensionReport[]} \\
projected trace & \texttt{Correlator[channel,"pert"]} \\
perturbative numerator & \texttt{PerturbativeNumerator[channel,s]} \\
spectral density & \texttt{PerturbativeSpectralDensity[channel,s]} \\
Borel moment & \texttt{NumericBorelPi[channel,order,M2,s0]} \\
mixing angle & \texttt{NumericMixingAngleDegrees[M2,s0,order]} \\
coordinate trace & \texttt{CoordinateTraceKernel[channel,order]} \\
coordinate angle & \texttt{CoordinateNumericMixingAngleDegrees[M2,s0,order]} \\
\end{tabular}
\end{center}

When a new term is added, the recommended debugging procedure is:
\begin{enumerate}
  \item compute one channel, usually \(AA\), at one point;
  \item check that the result is real;
  \item check its mass dimension;
  \item check its sign and size against the corresponding momentum-space
  expression, if available;
  \item only then scan over \(M^2\) and \(s_0\).
\end{enumerate}

\section{Correlators and spin-1 projection}

Define
\begin{equation}
  \Pi_{\mu\nu}^{ij}(p)
  =
  i\int d^4x\,e^{ip\cdot x}
  \langle 0 | T\{J^i_\mu(x)J^{j\dagger}_\nu(0)\}|0\rangle,
  \qquad i,j\in\{A,B\}.
\end{equation}

The invariant amplitude used for the axial-vector spin-1 state is extracted
with
\begin{equation}
  P^{\mu\nu}_{(1)}
  =
  \frac{1}{3}
  \left(
    g^{\mu\nu}
    -
    \frac{p^\mu p^\nu}{p^2}
  \right),
\end{equation}
so that
\begin{equation}
  \Pi_1^{ij}(p^2)=P^{\mu\nu}_{(1)}\Pi_{\mu\nu}^{ij}(p).
\end{equation}

The mixing angle is then obtained from
\begin{equation}
  \tan 2\theta
  =
  -\frac{2\Pi_1^{AB}}{\Pi_1^{AA}-\Pi_1^{BB}}.
\end{equation}
Numerically we use the quadrant-safe definition
\begin{equation}
  \theta_{\rm raw}
  =
  \frac{1}{2}
  \operatorname{atan2}
  \left(
    -2\Pi_1^{AB},
    \Pi_1^{AA}-\Pi_1^{BB}
  \right),
\end{equation}
and quote the principal branch
\begin{equation}
  \theta
  =
  \theta_{\rm raw}
  -
  \frac{\pi}{2}
  \operatorname{Round}
  \left(
    \frac{\theta_{\rm raw}}{\pi/2}
  \right),
  \qquad
  \theta\in[-45^\circ,45^\circ].
\end{equation}

\section{Momentum-space OPE}

The perturbative heavy-quark propagator is
\begin{equation}
  S_Q^{(0)}(k)
  =
  \frac{\slashed{k}+m_Q}{k^2-m_Q^2}.
\end{equation}

For the dimension-4 gluon condensate we use the background-field expansion
\begin{equation}
  S_Q(k)
  =
  S_Q^{(0)}(k)
  +
  g_sG_A^{\alpha\beta}T^A S_{Q,\alpha\beta}^{(G)}(k)
  +
  S_Q^{(G^2)}(k),
\end{equation}
where
\begin{align}
  S_{Q,\alpha\beta}^{(G)}(k)
  &=
  -\frac{1}{4}
  \frac{
    \sigma_{\alpha\beta}(\slashed{k}+m_Q)
    +
    (\slashed{k}+m_Q)\sigma_{\alpha\beta}
  }
  {(k^2-m_Q^2)^2},
  \\
  S_Q^{(G^2)}(k)
  &=
  \frac{\langle g_s^2G^2\rangle}{12}
  m_Q
  \frac{k^2+m_Q\slashed{k}}
  {(k^2-m_Q^2)^4}.
\end{align}
Equivalently,
\begin{equation}
  g_sG_A^{\alpha\beta}T^A S_{Q,\alpha\beta}^{(G)}(k)
  =
  -\frac{g_sG_A^{\alpha\beta}T^A}{4}
  \frac{
    \sigma_{\alpha\beta}(\slashed{k}+m_Q)
    +
    (\slashed{k}+m_Q)\sigma_{\alpha\beta}
  }
  {(k^2-m_Q^2)^2}.
\end{equation}
The one-gluon fields are contracted with
\begin{equation}
  \langle g_s^2G_A^{\alpha\beta}G_B^{\rho\sigma}\rangle
  =
  \frac{\delta_{AB}}{96}
  \langle g_s^2G^2\rangle
  \left(
    g^{\alpha\rho}g^{\beta\sigma}
    -
    g^{\alpha\sigma}g^{\beta\rho}
  \right),
\end{equation}
and
\begin{equation}
  \operatorname{Tr}(T^AT^B)=\frac{\delta^{AB}}{2}.
\end{equation}
Thus the color factor for the $S^{(G)}S^{(G)}$ piece is
\begin{equation}
  \frac{N_c^2-1}{192}\langle g_s^2G^2\rangle,
  \qquad N_c=3.
\end{equation}

The three projected OPE contributions are generated from
\begin{align}
  \Pi^{AA}_1 &\sim
  P^{\mu\nu}_{(1)}
  \operatorname{Tr}
  \left[
    \gamma_\mu\gamma_5
    S_c(k)
    \gamma_\nu\gamma_5
    S_b(k-p)
  \right],\\
  \Pi^{AB}_1 &\sim
  P^{\mu\nu}_{(1)}
  \operatorname{Tr}
  \left[
    \gamma_\mu\gamma_5
    S_c(k)
    i\sigma_{\nu\alpha}
    \frac{p^\alpha}{m_b+m_c}
    \gamma_5
    S_b(k-p)
  \right],\\
  \Pi^{BB}_1 &\sim
  P^{\mu\nu}_{(1)}
  \operatorname{Tr}
  \left[
    i\sigma_{\mu\alpha}
    \frac{p^\alpha}{m_b+m_c}
    \gamma_5
    S_c(k)
    i\sigma_{\nu\beta}
    \frac{p^\beta}{m_b+m_c}
    \gamma_5
    S_b(k-p)
  \right].
\end{align}

\section{Perturbative spectral densities}

For the perturbative loop, impose the two-particle cut
\begin{equation}
  k^2=m_c^2,\qquad (k-p)^2=m_b^2,
\end{equation}
which gives
\begin{equation}
  k\cdot p=\frac{s+m_c^2-m_b^2}{2},
  \qquad s=p^2.
\end{equation}

Define
\begin{equation}
  \lambda(s,m_b,m_c)
  =
  s^2+m_b^4+m_c^4
  -2sm_b^2-2sm_c^2-2m_b^2m_c^2 .
\end{equation}
The perturbative spectral densities are implemented as
\begin{equation}
  \rho_{\rm pert}^{ij}(s)
  =
  \frac{N_c}{16\pi^2}
  \frac{\sqrt{\lambda(s,m_b,m_c)}}{s}
  \frac{N^{ij}(s)}{(m_b+m_c)^{n_B^{ij}}},
\end{equation}
where $n_B^{AA}=0$, $n_B^{AB}=n_B^{BA}=1$, and $n_B^{BB}=2$ counts how many
tensor currents occur in the channel.  The unnormalized numerator functions
are
\begin{align}
  N^{AA}(s)
  &=
  -\frac{2}{s}
  \left[(m_b+m_c)^2-s\right]
  \left[(m_b-m_c)^2+2s\right],
  \\
  N^{AB}(s)
  &=
  6(m_b-m_c)\left[(m_b+m_c)^2-s\right],
  \\
  N^{BB}(s)
  &=
  2\left[
    -2(m_b^2-m_c^2)^2
    +(m_b^2-6m_bm_c+m_c^2)s
    +s^2
  \right].
\end{align}
Equivalently, after setting $N_c=3$,
\begin{align}
  \rho_{\rm pert}^{AA}(s)
  &=
  -\frac{3}{8\pi^2s^2}
  \left[(m_b+m_c)^2-s\right]
  \left[(m_b-m_c)^2+2s\right]
  \sqrt{\lambda(s,m_b,m_c)},\\
  \rho_{\rm pert}^{AB}(s)
  &=
  \frac{9}{8\pi^2s(m_b+m_c)}
  (m_b-m_c)
  \left[(m_b+m_c)^2-s\right]
  \sqrt{\lambda(s,m_b,m_c)},\\
  \rho_{\rm pert}^{BB}(s)
  &=
  \frac{3}{8\pi^2s(m_b+m_c)^2}
  \left[
    -2(m_b^2-m_c^2)^2
    +(m_b^2-6m_bm_c+m_c^2)s
    +s^2
  \right]
  \sqrt{\lambda(s,m_b,m_c)}.
\end{align}

The perturbative Borel moments are
\begin{equation}
  \Pi_{\rm pert}^{ij}(M^2,s_0)
  =
  \int_{(m_b+m_c)^2}^{s_0}
  ds\,
  e^{-s/M^2}
  \rho_{\rm pert}^{ij}(s).
\end{equation}

With this convention all three perturbative spectral densities have mass
dimension 2.  Without the factors of $(m_b+m_c)^{-n_B^{ij}}$, their
dimensions would be 2, 3, and 4 for $AA$, $AB$, and $BB$, respectively.

\section{Dimension-4 and dimension-6 gluon condensates}

The implemented condensate contribution is
\begin{equation}
  \Pi_{G^2}^{ij}
  =
  \Pi^{ij}[S_c^{(G^2)}S_b^{(0)}]
  +
  \Pi^{ij}[S_c^{(0)}S_b^{(G^2)}]
  +
  \Pi^{ij}[S_c^{(G)}S_b^{(G)}].
\end{equation}

After Feynman parametrization the terms have the schematic form
\begin{equation}
  \mathcal{A}_{G^2}^{ij}(x_1,x_2,s)
  =
  \sum_r
  C_r^{ij}(x_1,x_2,s)
  \left[
    m_c^2x_1^2
    +(m_b^2+m_c^2-s)x_1x_2
    +m_b^2x_2^2
  \right]^{-n_r}.
\end{equation}
The code then fixes the simplex by
\begin{equation}
  x_1=x,\qquad x_2=1-x.
\end{equation}
This gives
\begin{equation}
  \Delta(x,s)
  =
  m_c^2x+m_b^2(1-x)-s\,x(1-x).
\end{equation}

For the triple-gluon condensate we use the standard momentum-space
vacuum-averaged heavy-quark propagator term
\begin{equation}
  S_Q^{(G^3)}(k)
  =
  \frac{\langle g_s^3G^3\rangle}{48}
  \frac{
  (\slashed{k}+m_Q)
  \left[
    \slashed{k}(k^2-3m_Q^2)
    +2m_Q(2k^2-m_Q^2)
  \right]
  (\slashed{k}+m_Q)}
  {(k^2-m_Q^2)^6}.
\end{equation}
The implemented dimension-6 contribution is
\begin{equation}
  \Pi_{G^3}^{ij}
  =
  \Pi^{ij}[S_c^{(G^3)}S_b^{(0)}]
  +
  \Pi^{ij}[S_c^{(0)}S_b^{(G^3)}].
\end{equation}
This is the standard single-heavy-propagator contribution.  It does not yet
include any possible cross-line open-field terms of the schematic type
$S^{(G^2_{\rm open})}S^{(G)}$, which require an explicitly open two-gluon
propagator and a separate color/Lorentz derivation.

The $G^2$ and $G^3$ terms are not treated as ordinary smooth spectral
densities.  We instead use a direct Borel transform.  Set $s=-Q^2$ and write
\begin{equation}
  \Delta(x,-Q^2)
  =
  x(1-x)
  \left[Q^2+\bar s(x)\right],
\end{equation}
where
\begin{equation}
  \bar s(x)
  =
  \frac{m_c^2x+m_b^2(1-x)}{x(1-x)}.
\end{equation}

The Feynman-parameter amplitude is decomposed as
\begin{equation}
  -i\,\mathcal{A}_{G^r}^{ij}(x,-Q^2)
  =
  \sum_{\ell\geq 1}
  \frac{a_\ell^{ij}(x)}
  {\left[Q^2+\bar s(x)\right]^\ell}
  +
  \text{polynomial subtraction terms}.
\end{equation}
The polynomial subtraction terms vanish after the Borel transform.  The
working phase convention is the explicit factor $-i$ shown above.  This is a
bookkeeping convention for removing the loop-integration prefactor returned
by FeynCalc's Feynman-parameter amplitude.  This is the convention that
should be independently checked.

We use
\begin{equation}
  \mathcal{B}_{M^2}
  \left[
    \frac{1}{(Q^2+\bar s)^n}
  \right]
  =
  \frac{1}{(n-1)!\,(M^2)^{n-1}}
  e^{-\bar s/M^2}.
\end{equation}
Thus
\begin{equation}
  \Pi_{G^r}^{ij}(M^2,s_0)
  =
  \int_{x_-}^{x_+}
  dx\,
  e^{-\bar s(x)/M^2}
  \sum_{\ell\geq1}
  \frac{a_\ell^{ij}(x)}
  {(\ell-1)!\,(M^2)^{\ell-1}}.
\end{equation}

The continuum restriction is
\begin{equation}
  \bar s(x)\leq s_0,
\end{equation}
which gives
\begin{equation}
  x_\pm
  =
  \frac{
    s_0+m_b^2-m_c^2
    \pm
    \sqrt{\lambda(s_0,m_b,m_c)}
  }
  {2s_0}.
\end{equation}

\section{Total Borel moments and numerical test point}

The total moments used in the current numerical result are
\begin{equation}
  \Pi^{ij}(M^2,s_0)
  =
  \Pi_{\rm pert}^{ij}(M^2,s_0)
  +
  \Pi_{G^2}^{ij}(M^2,s_0)
  +
  \Pi_{G^3}^{ij}(M^2,s_0).
\end{equation}

The current default numerical inputs are
\begin{align}
  m_b &= 4.18~\text{GeV},\\
  m_c &= 1.27~\text{GeV},\\
  \langle g_s^2G^2\rangle &= 4\pi^2(0.012)~\text{GeV}^4,\\
  \langle g_s^3G^3\rangle &= 0.57~\text{GeV}^6,\\
  M^2 &= 10~\text{GeV}^2,\\
  s_0 &= 55~\text{GeV}^2.
\end{align}
For this point,
\begin{equation}
  x_- = 0.3322728664,\qquad
  x_+ = 0.9560816791.
\end{equation}

The perturbative moments are
\begin{align}
  \Pi_{\rm pert}^{AA} &= 0.1441716826,\\
  \Pi_{\rm pert}^{AB} &= -0.1051577192,\\
  \Pi_{\rm pert}^{BB} &= 0.1348134513.
\end{align}
The gluon-condensate moments are
\begin{align}
  \Pi_{G^2}^{AA} &= -0.0003658277,\\
  \Pi_{G^2}^{AB} &= 0.0002309420,\\
  \Pi_{G^2}^{BB} &= -0.0001888869.
\end{align}
The triple-gluon-condensate moments are
\begin{align}
  \Pi_{G^3}^{AA} &= 0.0000700979,\\
  \Pi_{G^3}^{AB} &= -0.0000493659,\\
  \Pi_{G^3}^{BB} &= 0.0000310238.
\end{align}
Thus
\begin{align}
  \Pi_{\rm total}^{AA} &= 0.1438759528,\\
  \Pi_{\rm total}^{AB} &= -0.1049761432,\\
  \Pi_{\rm total}^{BB} &= 0.1346555882.
\end{align}

The corresponding angles are
\begin{align}
  \theta_{\rm pert} &= 43.726118943^\circ,\\
  \theta_{\rm pert+G^2} &= 43.747426660^\circ,\\
  \theta_{\rm pert+G^2+G^3} &= 43.742693524^\circ.
\end{align}
At this point the relative condensate corrections are
\begin{align}
  \frac{\Pi_{G^2}^{AA}}{\Pi_{\rm pert}^{AA}} &= -0.002537,\\
  \frac{\Pi_{G^2}^{AB}}{\Pi_{\rm pert}^{AB}} &= -0.002196,\\
  \frac{\Pi_{G^2}^{BB}}{\Pi_{\rm pert}^{BB}} &= -0.001401,\\
  \frac{\Pi_{G^3}^{AA}}{\Pi_{\rm pert}^{AA}} &= 0.000486,\\
  \frac{\Pi_{G^3}^{AB}}{\Pi_{\rm pert}^{AB}} &= 0.000469,\\
  \frac{\Pi_{G^3}^{BB}}{\Pi_{\rm pert}^{BB}} &= 0.000230.
\end{align}

\section{Perturbative \texorpdfstring{$\alpha_s$}{alpha-s} sensitivity model}

The full perturbative $O(\alpha_s)$ correction is not yet calculated.  As a
diagnostic, the code implements a channel-dependent trial rescaling of the
leading perturbative moment,
\begin{equation}
  \Pi_{\rm pert}^{ij}(M^2,s_0)
  \longrightarrow
  \left[
    1+\frac{\alpha_s}{\pi}K_{ij}
  \right]
  \Pi_{\rm pert}^{ij}(M^2,s_0).
\end{equation}
The condensate pieces are left unchanged, so for the current total moment
\begin{equation}
  \Pi_{K}^{ij}
  =
  \left[
    1+\frac{\alpha_s}{\pi}K_{ij}
  \right]
  \Pi_{\rm pert}^{ij}
  +
  \Pi_{G^2}^{ij}
  +
  \Pi_{G^3}^{ij}.
\end{equation}
This is only a sensitivity model; it is not a substitute for the real
two-loop NLO perturbative spectral density.

For example, with $\alpha_s=0.26$ and
\begin{equation}
  K_{AA}=1,\qquad K_{AB}=0,\qquad K_{BB}=-1,
\end{equation}
the total angle at $M^2=10~\mathrm{GeV}^2$, $s_0=55~\mathrm{GeV}^2$ changes
from
\begin{equation}
  \theta_{\rm base}=43.742693524^\circ
\end{equation}
to
\begin{equation}
  \theta_K=40.625721673^\circ,
  \qquad
  \Delta\theta=-3.116971851^\circ.
\end{equation}
A coarse scan with
$K_{AA},K_{AB},K_{BB}\in\{-2,0,2\}$ gives a maximum branch-corrected shift of
about
\begin{equation}
  \max|\Delta\theta|\simeq 7.52^\circ.
\end{equation}

\section{Monte Carlo uncertainty propagation}

The numerical input table gives central values,
\begin{equation}
  m_b=4.18~\mathrm{GeV},\qquad
  m_c=1.27~\mathrm{GeV},
\end{equation}
\begin{equation}
  \langle g_s^2G^2\rangle = 4\pi^2(0.012)~\mathrm{GeV}^4,
  \qquad
  \langle g_s^3G^3\rangle = 0.57~\mathrm{GeV}^6.
\end{equation}
For the final quoted error these inputs should be assigned allowed intervals
or probability distributions.  In the present implementation we use uniform
sampling over editable intervals.  A Monte Carlo point is
\begin{equation}
  X_i =
  \left\{
    m_b^{(i)},m_c^{(i)},G_2^{(i)},G_3^{(i)},
    M_i^2,s_{0,i}
  \right\},
\end{equation}
where
\begin{equation}
  G_2\equiv\langle g_s^2G^2\rangle,\qquad
  G_3\equiv\langle g_s^3G^3\rangle.
\end{equation}
Each component is drawn from the user-specified interval
\begin{equation}
  x^{(i)}\sim U(x_{\min},x_{\max}).
\end{equation}
Fixed parameters are implemented by setting \(x_{\min}=x_{\max}\), or in the
Mathematica code by passing a single number instead of an interval.

The current working ranges in the code are
\begin{align}
  m_b &\in [4.16,4.21]~\mathrm{GeV},&
  m_c &\in [1.25,1.29]~\mathrm{GeV},\\
  G_2 &\in 4\pi^2[0.006,0.018]~\mathrm{GeV}^4,&
  G_3 &\in [0.28,0.86]~\mathrm{GeV}^6,\\
  M^2 &\in [8,12]~\mathrm{GeV}^2,&
  s_0 &\in [50,60]~\mathrm{GeV}^2.
\end{align}
Equivalently, in the central-value-plus-error notation useful for the paper,
\begin{align}
  m_b &= (4.185\pm0.025)~\mathrm{GeV},&
  m_c &= (1.27\pm0.02)~\mathrm{GeV},\\
  G_2 &= 4\pi^2(0.012\pm0.006)~\mathrm{GeV}^4,&
  G_3 &= (0.57\pm0.29)~\mathrm{GeV}^6,\\
  M^2 &= (10\pm2)~\mathrm{GeV}^2,&
  s_0 &= (55\pm5)~\mathrm{GeV}^2.
\end{align}
These are not final physics choices. They are placeholders for the scan layer
and should be replaced by the intervals adopted in the analysis.

For every random point we impose the physical continuum condition
\begin{equation}
  s_{0,i} > \left(m_b^{(i)}+m_c^{(i)}\right)^2.
\end{equation}
Accepted points are evaluated with the same OPE formula as above,
\begin{equation}
  \theta_i =
  \frac{1}{2}\operatorname{ArcTan}
  \left[
    \Pi_1^{AA}(X_i)-\Pi_1^{BB}(X_i),
    -2\Pi_1^{AB}(X_i)
  \right],
\end{equation}
using the chosen OPE order, normally ``total'' for
\(\Pi_{\rm pert}+\Pi_{G^2}+\Pi_{G^3}\).
For the present fast uncertainty scan we keep the triple-gluon term switched
off by default.  In the Mathematica implementation this is controlled by
\texttt{\$BcMixingMonteCarloIncludeG3=False} or by the option
\texttt{"IncludeG3" -> False}.  With this switch, a requested ``total'' Monte
Carlo scan is evaluated as ``pertG2'',
\begin{equation}
  \Pi_{\rm MC}^{ij}
  =
  \Pi_{\rm pert}^{ij}+\Pi_{G^2}^{ij},
\end{equation}
and the \(G^3\) contribution can be restored later by setting
\texttt{"IncludeG3" -> True}.

The central value and Gaussian-width estimate are then
\begin{equation}
  \bar\theta = \frac{1}{N}\sum_{i=1}^N \theta_i,
  \qquad
  \sigma_{\theta,\mathrm{Gauss}} =
  \sqrt{\frac{1}{N}\sum_{i=1}^N
  \left(\theta_i-\bar\theta\right)^2}.
\end{equation}
The code also reports the unbiased sample standard deviation, median, minimum,
maximum, and the \(16\%\) and \(84\%\) quantiles. These quantiles are useful
because the resulting \(\theta\) distribution may not be exactly Gaussian.

The Mathematica entry point is
\begin{verbatim}
nSamples = 500;
histBins = 25;
m2Range = {8.0, 12.0};
s0Range = {50.0, 60.0};

ranges = Join[
  $BcMixingDefaultUncertaintyRanges,
  <|"M2" -> m2Range, "s0" -> s0Range|>
];

uncertaintyRun = MonteCarloMixingAngleUncertainty[
  nSamples, ranges, "total",
  "IncludeG3" -> False,
  "Seed" -> 1234,
  "Progress" -> True
];
uncertaintyRun["Summary"]
MonteCarloMixingAngleHistogram[uncertaintyRun, histBins]
MonteCarloMixingAnglePublicationHistogram[uncertaintyRun, histBins]
\end{verbatim}
where \texttt{ranges} is an association such as
\begin{verbatim}
<|
  "mb" -> {4.16, 4.21},
  "mc" -> {1.25, 1.29},
  "G2" -> {4 Pi^2 0.006, 4 Pi^2 0.018},
  "G3" -> {0.28, 0.86},
  "M2" -> {8.0, 12.0},
  "s0" -> {50.0, 60.0}
|>
\end{verbatim}
The number of random points, the input intervals, and the Borel window are
therefore all editable without touching the algebra.

The first argument, for example \texttt{500}, is the number of random accepted
Monte Carlo points. The option \texttt{"Seed" -> 1234} fixes the random-number
seed so the sample is reproducible. The Monte Carlo result is called
\texttt{uncertaintyRun}, not \texttt{mc}, to avoid confusion with the charm
mass parameter \texttt{"mc"}. To vary only the Borel window and the
continuum threshold while keeping the QCD inputs fixed at their central values,
one can use
\begin{verbatim}
ranges = Join[
  $BcMixingDefaultUncertaintyRanges,
  <|
    "mb" -> 4.18,
    "mc" -> 1.27,
    "G2" -> 4 Pi^2 0.012,
    "G3" -> 0.57,
    "M2" -> {8.0, 12.0},
    "s0" -> {50.0, 60.0}
  |>
];
\end{verbatim}
The accepted random points and the summary can be exported for replotting:
\begin{verbatim}
ExportMonteCarloMixingAngleSamples[uncertaintyRun, "BcMixingMonteCarloSamples.csv"];
ExportMonteCarloMixingAngleSummary[uncertaintyRun, "BcMixingMonteCarloSummary.csv"];
Export["BcMixingMonteCarloHistogram.pdf",
  MonteCarloMixingAnglePublicationHistogram[uncertaintyRun, 25]];
\end{verbatim}

\section{How to reproduce the numerical check}

In Mathematica, load the implementation with
\begin{verbatim}
SetDirectory["/Users/sbilmis/Bc_mixing"];
Get["BcMixingMomentum.wl"];
\end{verbatim}
Then evaluate
\begin{verbatim}
NumericOPESummary[10, 55]
NumericMixingAngleDegrees[10, 55, "pert"]
NumericMixingAngleDegrees[10, 55, "pertG2"]
NumericMixingAngleDegrees[10, 55, "total"]
KFactorSensitivityRecord[10, 55, KFactorAssociation[1, 0, -1], 0.26, "total"]
KFactorSensitivityEnvelope[10, 55, {-2, 2, 2}, 0.26, "total"]
MixingMassDimensionReport[]
CheckMixingMassDimensions[]
PerturbativeSpectralDensityDimensionReport[]
\end{verbatim}
For window studies one can scan
\begin{verbatim}
MixingAngleTable[{8, 12, 1}, {50, 60, 5}, "total"]
OPEConvergenceDataset[{8, 12, 1}, {50, 60, 5}]
uncertaintyRun = MonteCarloMixingAngleUncertainty[50, $BcMixingDefaultUncertaintyRanges,
  "total", "IncludeG3" -> False, "Seed" -> 1234]
uncertaintyRun["Summary"]
\end{verbatim}

\section{Auxiliary-parameter stability plots}

As a comparison with the vector \(B_c^*\) analysis of Wang, we also consider
the auxiliary window
\begin{equation}
  M^2\in[7.0,9.0]~\mathrm{GeV}^2,
  \qquad
  s_0=(54\pm1)~\mathrm{GeV}^2.
\end{equation}
The Mathematica implementation stores this as
\begin{verbatim}
$BcMixingWangWindow
\end{verbatim}
and produces publication-style stability plots with
\begin{verbatim}
stabilityOrder = "pertG2";

m2Stability = MixingAngleStabilityM2PublicationPlot[
  $BcMixingWangWindow["M2Range"],
  $BcMixingWangWindow["s0Values"],
  stabilityOrder,
  $BcMixingDefaultParameters,
  "NPoints" -> 25,
  "YHalfWidth" -> 1.0
];

s0Stability = MixingAngleStabilityS0PublicationPlot[
  $BcMixingWangWindow["s0Range"],
  $BcMixingWangWindow["M2Values"],
  stabilityOrder,
  $BcMixingDefaultParameters,
  "NPoints" -> 25,
  "YHalfWidth" -> 1.0
];
\end{verbatim}
The option \texttt{"YHalfWidth" -> 1.0} prevents the vertical axis from being
over-zoomed; it enforces at least a \(\pm1^\circ\) display window around the
calculated angles.  This is useful when the goal is to demonstrate stability
with respect to the auxiliary parameters, not to magnify tiny numerical
variations.
For publication labels the plotting functions try to use the Mathematica
package MaTeX automatically,
\begin{verbatim}
"UseMaTeX" -> Automatic
\end{verbatim}
so that the axes and legends are rendered as \(M^2(\mathrm{GeV}^2)\),
\(s_0(\mathrm{GeV}^2)\), \(\theta^\circ\), and entries such as
\(s_0=54\,\mathrm{GeV}^2\).  If MaTeX is not installed, the code falls back to
ordinary Mathematica labels with a degree symbol.  The fallback can be forced
with \texttt{"UseMaTeX" -> False}.

The plots and their underlying data can be exported with
\begin{verbatim}
Export["BcMixingThetaVsM2_WangWindow.pdf", m2Stability["Plot"]];
Export["BcMixingThetaVsS0_WangWindow.pdf", s0Stability["Plot"]];
ExportStabilityDataCSV[m2Stability["Data"],
  "BcMixingThetaVsM2_WangWindow.csv", "M2"];
ExportStabilityDataCSV[s0Stability["Data"],
  "BcMixingThetaVsS0_WangWindow.csv", "s0"];
\end{verbatim}

\section{Coordinate-space cross-check}

The file \texttt{BcMixingCoordinate.wl} starts an independent coordinate-space
validation of the perturbative part.  This is meant as a blind check of the
momentum-space calculation, not as a replacement of it.  The same interpolating
currents and the same normalization of the tensor current are used,
\begin{equation}
  J^A_\mu(x)=\bar b(x)\gamma_\mu\gamma_5c(x),
  \qquad
  J^B_\mu(x)=
  i\bar b(x)\sigma_{\mu\alpha}
  \frac{p^\alpha}{m_b+m_c}\gamma_5c(x).
\end{equation}
With these conventions, the coordinate-space perturbative correlators are
\begin{align}
  \Pi_{\mu\nu}^{AA}(p)
  &=
  -iN_c\int d^4x\,e^{ip\cdot x}
  \mathrm{Tr}\!\left[
    \gamma_\mu\gamma_5 S_c^{(0)}(-x)
    \gamma_\nu\gamma_5 S_b^{(0)}(x)
  \right],
  \\
  \Pi_{\mu\nu}^{AB}(p)
  &=
  -iN_c\int d^4x\,e^{ip\cdot x}
  \mathrm{Tr}\!\left[
    \gamma_\mu\gamma_5 S_c^{(0)}(-x)
    i\sigma_{\nu\beta}\frac{p^\beta}{m_b+m_c}\gamma_5
    S_b^{(0)}(x)
  \right],
  \\
  \Pi_{\mu\nu}^{BB}(p)
  &=
  -iN_c\int d^4x\,e^{ip\cdot x}
  \mathrm{Tr}\!\left[
    i\sigma_{\mu\alpha}\frac{p^\alpha}{m_b+m_c}\gamma_5
    S_c^{(0)}(-x)
    i\sigma_{\nu\beta}\frac{p^\beta}{m_b+m_c}\gamma_5
    S_b^{(0)}(x)
  \right].
\end{align}
The sign difference between \(S_c(-x)\) and \(S_b(x)\) is implemented in the
code by changing the sign of the \(\slashed{x}\) term.  The spin--1 invariant
amplitude is extracted with the same transverse projector as in momentum
space,
\begin{equation}
  \Pi_1^{ij}(p^2)
  =
  \frac{1}{3}
  \left(g^{\mu\nu}-\frac{p^\mu p^\nu}{p^2}\right)
  \Pi_{\mu\nu}^{ij}(p).
\end{equation}
This step is performed in \texttt{CoordinateTraceKernel[channel,"pert"]}.

The free heavy-quark propagator is written in terms of modified Bessel
functions,
\begin{equation}
  S_Q^{(0)}(x)
  =
  \frac{m_Q^2}{(2\pi)^2}
  \left[
    i\slashed{x}
    \frac{K_2(m_Q\sqrt{-x^2})}{-x^2}
    +
    \frac{K_1(m_Q\sqrt{-x^2})}{\sqrt{-x^2}}
  \right].
\end{equation}
This is the form used in the coordinate-space treatments of heavy systems.
The derivative-free perturbative trace contains products of the form
\begin{equation}
  (\slashed{x})^a (p\cdot x)^b
  \frac{K_\nu(m_c\sqrt{-x^2})K_\mu(m_b\sqrt{-x^2})}
       {(\sqrt{-x^2})^n},
\end{equation}
where the tensor numerator is reduced by FeynCalc before the radial Fourier
integral is performed.

After Wick rotation to Euclidean momentum \(Q^2=-p^2\), the angular integral
has the standard form
\begin{equation}
  i\int d^Dx\,e^{ip\cdot x}f(\sqrt{-x^2})
  =
  (2\pi)^{D/2}Q^{1-D/2}
  \int_0^\infty dx\,x^{D/2}J_{D/2-1}(Qx)f(x).
\end{equation}
Therefore every scalar term is reduced to radial integrals
\begin{equation}
  I_{u\nu\mu}(Q;m_c,m_b)
  =
  \int_0^\infty dx\,x^{u-1}
  J_v(Qx)K_\nu(m_cx)K_\mu(m_bx).
\end{equation}
For equal heavy masses, Huang--Liu give a compact hypergeometric expression
for this integral.  For two different heavy masses, the Azizi et al.
Bessel/Schwinger representation leads to the same physical two-body cut.  The
installed perturbative spectral density is
\begin{equation}
  \rho_{\rm pert}^{ij}(s)
  =
  \frac{N_c}{16\pi^2}
  \frac{\sqrt{\lambda(s,m_b,m_c)}}{s}\,
  N^{ij}(s)\,
  \left(m_b+m_c\right)^{-n_{ij}},
\end{equation}
where \(n_{AA}=0\), \(n_{AB}=1\), \(n_{BB}=2\), and
\begin{align}
  \lambda(s,m_b,m_c)
  &=
  s^2+m_b^4+m_c^4
  -2sm_b^2-2sm_c^2-2m_b^2m_c^2,\\
  k\cdot p\big|_{\rm cut}
  &=
  \frac{s+m_c^2-m_b^2}{2}.
\end{align}
The numerator functions are
\begin{align}
  N^{AA}(s)
  &=
  -\frac{4}{s}
  \left[
    2(k\cdot p)^2
    +(3m_bm_c+m_c^2)s
    -3s(k\cdot p)
  \right],
  \\
  N^{AB}(s)
  &=
  12\left[-(m_b+m_c)(k\cdot p)+m_cs\right],
  \\
  N^{BB}(s)
  &=
  -4
  \left[
    4(k\cdot p)^2
    +(3m_bm_c-m_c^2)s
    -3s(k\cdot p)
  \right].
\end{align}
A hand check should verify these three numerators from the projected
coordinate-space trace and then verify the phase-space factor
\(N_c\sqrt{\lambda}/(16\pi^2s)\).  In the code these expressions are stored in
\texttt{CoordinatePerturbativeNumerator} and
\texttt{CoordinatePerturbativeSpectralDensity}.

The perturbative Borel moment is then evaluated as
\begin{equation}
  \Pi_{\rm pert}^{ij}(M^2,s_0)
  =
  \int_{(m_b+m_c)^2}^{s_0}ds\,
  e^{-s/M^2}\rho_{\rm pert}^{ij}(s),
\end{equation}
and the angle is obtained from the same diagonalization formula,
\begin{equation}
  \theta_{\rm coord}^{\rm pert}(M^2,s_0)
  =
  \frac12
  \operatorname{ArcTan}
  \left[
    \Pi_{\rm pert}^{AA}-\Pi_{\rm pert}^{BB},
    -2\Pi_{\rm pert}^{AB}
  \right].
\end{equation}

\subsection{What should be written in the journal paper}

If the paper presents the coordinate-space calculation as the main method,
the main text should not merely state that the result was obtained with
Mathematica.  It should show enough of the analytic reduction that a reader
can see where the physical cut and continuum subtraction enter.  A compact
but sufficient presentation would include:
\begin{enumerate}
  \item the two currents, including the tensor-current normalization
  \(p^\alpha/(m_b+m_c)\);
  \item the coordinate-space correlators after Wick contraction;
  \item the heavy-quark propagator in terms of \(K_\nu(m_Q\sqrt{-x^2})\);
  \item the spin--1 projection used to define \(\Pi_1^{ij}\);
  \item the generic Bessel integral generated by the Fourier transform;
  \item the way the Bessel functions are converted to Schwinger-parameter
  integrals;
  \item the emergence of Dirac delta functions after the Borel transform;
  \item the resulting Heaviside support conditions or, equivalently, the
  integration limits that implement continuum subtraction;
  \item the final perturbative spectral densities, preferably in the compact
  \(N^{ij}(s)\sqrt{\lambda}/s\) form;
  \item the statement that the current coordinate-space condensate check
  includes local \(G^2\) and \(G^3\) single-line terms, while the full
  dimension--4 coordinate-space result additionally needs the
  \(S_c^{(G)}S_b^{(G)}\) cross-line term.
\end{enumerate}
The main text can contain the compact formulas and the final spectral
densities.  Longer channel-by-channel trace expressions and the local
condensate Borel kernels are better placed in an appendix or supplemental
file.

\subsection{Role of the Azizi et al. method}

The main technical role of the Azizi et al. method is to make continuum
subtraction well defined for systems with more than one heavy propagator in
coordinate space.  With one heavy quark, the coordinate-space Fourier
transform contains one modified Bessel function and can often be reduced to
ordinary hypergeometric functions.  With two heavy quarks, the Fourier
transform contains products
\begin{equation}
  J_v(Qx)K_\nu(m_cx)K_\mu(m_bx),
\end{equation}
and a naive treatment of the Bessel functions can lead to parameter integrals
with apparent endpoint divergences or ambiguous continuum subtraction.

Azizi et al. avoid this by using an integral representation of the modified
Bessel functions and then applying the Borel transformations before doing all
parameter integrations.  In this order of operations, the Borel transforms
generate Dirac delta functions.  These delta functions fix combinations of
Schwinger parameters and convert the would-be divergent parameter domains into
finite domains with explicit support conditions.  The support conditions are
usually written as Heaviside functions or, after solving the inequalities, as
modified integration limits.

Without this organization one would have to handle the delta functions,
derivatives of delta functions, endpoint constraints, and continuum subtraction
analytically by hand for each term.  That is possible, but it is error-prone
for two heavy masses because every numerator power changes the derivative
order of the delta function or the power of the support factor.  The Azizi
formulation gives a systematic recipe:
\begin{align}
  K_\nu K_\mu
  &\longrightarrow
  \text{Schwinger parameters}
  \longrightarrow
  \text{Borel delta functions}
  \nonumber\\
  &\longrightarrow
  \Theta\text{-functions}
  \longrightarrow
  \text{finite integration limits}.
\end{align}
In the paper, this chain should be shown once for a generic term.  The final
channel-dependent spectral densities can then be quoted in compact form.

\subsection{Direct coordinate-space route without the Azizi reduction}

It is useful to state explicitly what would happen if one deliberately did
not use the Azizi et al. organization.  This is not the route used by the
current numerical code, and it should be treated as an independent analytic
cross-check or appendix-level exercise.  The starting point is still the same
coordinate-space correlator,
\begin{equation}
  \Pi_{\mu\nu}^{ij}(p)
  =
  -iN_c\int d^4x\,e^{ip\cdot x}
  \mathrm{Tr}\!\left[
    \Gamma_\mu^i S_c(-x)\Gamma_\nu^j S_b(x)
  \right],
\end{equation}
but after the Dirac trace one keeps the Fourier transform of each Bessel
product directly.  A typical scalar integral has the form
\begin{equation}
  \mathcal I_{n\mu\nu}^{(r)}(p^2)
  =
  \int d^4x\,e^{ip\cdot x}
  (p\cdot x)^r(x^2)^n
  K_\mu(m_c\sqrt{-x^2})K_\nu(m_b\sqrt{-x^2}) .
\end{equation}
The powers of \(p\cdot x\) may be generated by differentiating the scalar
Fourier integral with respect to \(p_\alpha\),
\begin{equation}
  x_\alpha e^{ip\cdot x}
  =
  \frac{1}{i}\frac{\partial}{\partial p^\alpha}e^{ip\cdot x},
\end{equation}
or by doing Lorentz tensor reduction before the radial integral.  The angular
integration then gives integrals containing
\begin{equation}
  J_1(Qx)K_\mu(m_cx)K_\nu(m_bx),
  \qquad Q^2=-p^2 .
\end{equation}
At this stage the direct method must extract the discontinuity in
\(p^2=s\) from the Bessel product itself.  If Schwinger parameters are
introduced only as a calculational aid, one meets distributions of the type
\begin{equation}
  \delta\!\left(s-\frac{m_c^2}{x}
                 -\frac{m_b^2}{1-x}\right),
\end{equation}
and, when numerator powers or derivatives act on the Fourier kernel,
derivatives of the same delta function,
\begin{equation}
  \delta^{(n)}\!\left(s-\bar s(x)\right),
  \qquad
  \bar s(x)=\frac{m_c^2}{x}+\frac{m_b^2}{1-x}.
\end{equation}
These objects must be treated analytically.  For example,
\begin{equation}
  \int_0^1 dx\,F(x)\,
  \delta\!\left(s-\bar s(x)\right)
  =
  \sum_{x_i}
  \frac{F(x_i)}
       {|\bar s'(x_i)|},
  \qquad
  \bar s(x_i)=s,
\end{equation}
where the roots \(x_i\) exist only if \(s\geq(m_b+m_c)^2\).  This condition
is the same physical threshold and may be written as a Heaviside factor,
\begin{equation}
  \Theta\!\left(s-(m_b+m_c)^2\right).
\end{equation}
For derivative delta functions one must integrate by parts,
\begin{equation}
  \int dx\,F(x)\,\delta^{(n)}(s-\bar s(x)),
\end{equation}
carefully keeping the Jacobian factors and verifying that no endpoint terms
survive after the physical support restriction is imposed.

Thus the direct route has the schematic structure
\begin{align}
  \text{Bessel product}
  &\longrightarrow
  \text{radial Fourier transform}
  \longrightarrow
  \operatorname{Disc}_{p^2=s}
  \nonumber\\
  &\longrightarrow
  \delta(s-\bar s),\ \delta'(s-\bar s),\ldots
  \longrightarrow
  \Theta(s-s_{\rm th})
  \longrightarrow
  x_\pm(s).
\end{align}
This is analytically possible, but it is precisely the bookkeeping that the
Azizi et al. method systematizes.  The direct section is therefore best used
as a validation path: derive one channel, preferably \(AA\), all the way to
the same perturbative spectral density and confirm that the final answer is
proportional to
\begin{equation}
  \frac{\sqrt{\lambda(s,m_b,m_c)}}{s}
  \Theta\!\left(s-(m_b+m_c)^2\right).
\end{equation}
Only after the \(AA\) channel is reproduced should one repeat the calculation
for \(AB\) and \(BB\), because those channels additionally test the tensor
current normalization \(1/(m_b+m_c)\).  The condensate terms are substantially
more cumbersome in this route, since local \(G^2\), local \(G^3\), and the
open-field \(S_c^{(G)}S_b^{(G)}\) term can generate higher derivatives of the
same support delta functions.

\subsection{Where the Dirac delta and Heaviside functions enter}

The reason the current Mathematica output does not visibly display
\(\delta\)-functions and Heaviside functions is not that they are absent from
the coordinate-space method.  Rather, in the current implementation they have
already been integrated out or replaced by equivalent integration limits.

The coordinate-space product of two massive propagators generates scalar
terms of the schematic form
\begin{equation}
  T
  =
  \int d^4x\,e^{ip\cdot x}
  \frac{
    K_\nu(m_c\sqrt{-x^2})K_\mu(m_b\sqrt{-x^2})
  }{(\sqrt{-x^2})^n}
  \times \{\text{polynomial in }x^2,p\cdot x,p^2\}.
\end{equation}
Using an integral representation for each modified Bessel function, one obtains
Schwinger-parameter integrals of the schematic form
\begin{equation}
  T
  \sim
  \int d\rho\,dv\,dw\,
  \rho^a v^b w^c(1-v-w)^d
  \exp\left[
    -\frac{m_c^2}{\rho v}
    -\frac{m_b^2}{\rho w}
    -\frac{Q^2}{4\rho}
  \right],
  \qquad Q^2=-p^2 .
\end{equation}
This is the stage emphasized in the Azizi et al. treatment.  In applications
with two external virtualities, successive Borel transformations produce
Dirac delta functions that fix combinations of the Schwinger parameters.  A
typical reduced structure is
\begin{equation}
  \rho(s_1,s_2)
  \sim
  \int dz\,dy\,dl\,
  F(z,y,l)\,
  \delta\!\left[
    s_1-l-\frac{m_c^2}{z(1-y)}-\frac{m_b^2}{zy}
  \right]
  \delta(s_2-s_1),
\end{equation}
up to derivatives of the delta functions when powers of the external
invariants appear.  After the \(l\)-integration, these delta functions turn
into support restrictions,
\begin{align}
  &\int_0^\infty dl\,
  F(l)\,
  \delta\!\left(
    s-l-A
  \right)
  =
  F(s-A)\,\Theta(s-A),
  \\
  &A
  =
  \frac{m_c^2}{z(1-y)}
  +
  \frac{m_b^2}{zy}.
\end{align}
Thus one obtains
\begin{equation}
  \Theta\!\left(
    s_1-\frac{m_c^2}{z(1-y)}-\frac{m_b^2}{zy}
  \right),
\end{equation}
or equivalently into restricted integration limits.  These \(\Theta\)-functions
are what remove apparent endpoint divergences in the Schwinger-parameter
integrals.

For the present two-point mixing correlator there is one external invariant,
\(p^2=s\), and the same physics appears in a simpler form: the spectral density
has support only above
\begin{equation}
  s_{\rm th}=(m_b+m_c)^2,
\end{equation}
or, in a Feynman-parameter representation,
\begin{equation}
  s >
  \bar s(x)
  =
  \frac{m_c^2x+m_b^2(1-x)}{x(1-x)} .
\end{equation}
Solving \(\bar s(x)<s_0\) gives the finite integration interval
\begin{equation}
  x_-(s_0)<x<x_+(s_0),
\end{equation}
where
\begin{equation}
  x_\pm(s_0)
  =
  \frac{s_0+m_b^2-m_c^2
  \pm\sqrt{\lambda(s_0,m_b,m_c)}}{2s_0}.
\end{equation}
This is the explicit replacement used in the numerical integrations:
\begin{equation}
  \int_0^1 dx\,
  F(x)\,
  \Theta\!\left(s_0-\bar s(x)\right)
  \quad\longrightarrow\quad
  \int_{x_-(s_0)}^{x_+(s_0)}dx\,F(x).
\end{equation}
Equivalently, when writing the dispersion integral, the same support appears
as
\begin{equation}
  \int ds\,\rho(s)
  =
  \int ds\,
  \Theta\!\left(s-(m_b+m_c)^2\right)
  \rho(s),
\end{equation}
with the continuum-subtracted Borel moment
\begin{equation}
  \Pi(M^2,s_0)
  =
  \int_{(m_b+m_c)^2}^{s_0}
  ds\,e^{-s/M^2}\rho(s).
\end{equation}
Thus the Heaviside support has not disappeared; it is encoded in the lower
threshold and in the \(x_\pm\) integration limits used by the numerical Borel
integrals.

The present coordinate-space code has two layers:
\begin{enumerate}
  \item For the perturbative check, it installs the final reduced spectral
  density proportional to
  \(\sqrt{\lambda(s,m_b,m_c)}\,N^{ij}(s)/s\).  Therefore no intermediate
  \(\delta\)- or \(\Theta\)-functions appear in the output.
  \item For the local \(G^2\) and \(G^3\) checks, it uses the already-tested
  Feynman/Schwinger direct-Borel reduction helpers from
  \texttt{BcMixingMomentum.wl}.  In this representation the support is again
  implemented by the integration limits \(x_\pm(s_0)\), not by printing
  explicit Heaviside functions.
\end{enumerate}
For a journal paper, it is therefore important to say explicitly that the
Dirac delta functions and Heaviside functions occur at the Schwinger/Borel
reduction stage, but in the numerical implementation they are already
integrated out and appear as physical thresholds and finite parameter
integration domains.

For a student checking this section by hand, the suggested coordinate-space
exercise is:
\begin{enumerate}
  \item Start from the \(AA\) current only.  Insert \(S_c^{(0)}(-x)\) and
  \(S_b^{(0)}(x)\) and keep track of the sign of the \(\slashed{x}\) term.
  \item Evaluate the Dirac trace before projection.  The trace should reduce
  to scalar combinations of \(x^2\), \(p^2\), and \(p\cdot x\), multiplied by
  products of \(K_1\) and \(K_2\).
  \item Apply the transverse projector and replace
  \(p^2\to s\), \(x^2\to x^2\), and \(p\cdot x\to px\).  This is what
  \texttt{CoordinateScalarize} does.
  \item Identify the radial integral associated with each scalar term.  Terms
  with powers of \(p\cdot x\) can be generated by differentiating the Fourier
  factor with respect to \(p_\mu\), or equivalently by using tensor reduction
  before the angular integral.
  \item Use the Bessel/Schwinger representation to obtain the physical cut in
  \(s\).  The final perturbative answer must be proportional to
  \(\sqrt{\lambda(s,m_b,m_c)}\) and vanish at \(s=(m_b+m_c)^2\).
  \item Repeat the numerator check for \(AB\) and \(BB\).  These two channels
  are the best tests of the tensor-current normalization because each \(B\)
  current contributes one factor \(1/(m_b+m_c)\).
\end{enumerate}

The coordinate-space notebook separates these checks intentionally:
\begin{itemize}
  \item first evaluate only \texttt{BcMixingCoordinate.wl};
  \item record \texttt{CoordinateTraceKernel["AA","pert"]};
  \item record \texttt{CoordinatePerturbativeSpectralDensity["AA",s]};
  \item compute \texttt{CoordinateNumericMixingAngleDegrees[8,54,"pert"]};
  \item only after this blind check, load \texttt{BcMixingMomentum.wl} and
  compare with
\begin{verbatim}
CoordinateCompareToMomentum[8,54]
\end{verbatim}
\end{itemize}

The Mathematica entry points are
\begin{verbatim}
Get["BcMixingCoordinate.wl"];
CoordinateCheckEnvironment[]
CoordinateTraceKernel["AA", "pert"] // Short
CoordinateNumericMixingAngleDegrees[8, 54, "pert"]
CoordinateWangWindowGrid[]
\end{verbatim}
If the momentum-space file is also loaded, the two perturbative calculations
can be compared by
\begin{verbatim}
Get["BcMixingMomentum.wl"];
Get["BcMixingCoordinate.wl"];
CoordinateCompareToMomentum[8, 54]
\end{verbatim}
which gives
\begin{verbatim}
<|"CoordinatePertDeg" -> 43.28944792785432,
  "MomentumPertDeg" -> 43.28944792785432,
  "DifferenceDeg" -> 0.|>
\end{verbatim}
for \(M^2=8~\mathrm{GeV}^2\) and \(s_0=54~\mathrm{GeV}^2\).

For the next stage, the coordinate-space \(G^2\) and \(G^3\) propagator
kernels are already written as extension points.  The local-condensate
heavy-quark terms used in the code are
\begin{align}
  S_Q^{(G^2)}(x)
  &=
  -\frac{\langle g_s^2G^2\rangle}
  {2^6\,3\,(2\pi)^2}
  \left[
    (im_Q\slashed{x}-6)
    \frac{K_1(m_Q\sqrt{-x^2})}{\sqrt{-x^2}}
    +4m_Qx^2
    \frac{K_2(m_Q\sqrt{-x^2})}{-x^2}
  \right],
  \\
  S_Q^{(G^3)}(x)
  &=
  \frac{\langle g_s^3G^3\rangle}
  {2^8\,3^2\,(2\pi)^2}
  \left[
    -\frac{i\slashed{x}x^2}{m_Q}
    \frac{K_1(m_Q\sqrt{-x^2})}{\sqrt{-x^2}}
    +4i\slashed{x}x^2
    \frac{K_2(m_Q\sqrt{-x^2})}{-x^2}
  \right.
  \nonumber\\
  &\hspace{4.0cm}\left.
    +\frac{10x^4}{m_Q}
    \frac{K_2(m_Q\sqrt{-x^2})}{-x^2}
    +x^4
    \frac{K_1(m_Q\sqrt{-x^2})}{\sqrt{-x^2}}
  \right].
\end{align}
The dimension--4 coordinate-space check must include
\begin{equation}
  S_c^{(G^2)}(-x)S_b^{(0)}(x),\qquad
  S_c^{(0)}(-x)S_b^{(G^2)}(x),
\end{equation}
and, if the open-field propagator is kept explicitly, the cross-line
\(S_c^{(G)}(-x)S_b^{(G)}(x)\) term.  The present coordinate file records the
local \(G^2\) and \(G^3\) kernels and can build their projected trace kernels,
for example
\begin{verbatim}
CoordinateTraceKernel["AA", "G2c"]      (* S_c^(G2) S_b^(0) *)
CoordinateTraceKernel["AA", "G2b"]      (* S_c^(0)  S_b^(G2) *)
CoordinateTraceKernel["AA", "G2local"]  (* local sum of the two *)
CoordinateTraceKernel["AA", "G3local"]  (* local single-line G3 sum *)
\end{verbatim}
At this stage, however, the numerical coordinate-space spectral/Borel
integration is enabled directly only for \texttt{"pert"}.  A reduced
Bessel/Schwinger integration bridge has now been added for the local
single-line condensate pieces.  It is used only after loading
\texttt{BcMixingMomentum.wl}, because that file already contains the tested
Feynman/Schwinger parameter reduction and direct Borel transform.  The
coordinate notebook therefore keeps these cells in an optional section:
\begin{verbatim}
Get["BcMixingMomentum.wl"];
Get["BcMixingCoordinate.wl"];

CoordinateLocalCondensateSummary[8, 54]
CoordinateNumericMixingAngleDegrees[8, 54, "pertG2local"]
CoordinateNumericMixingAngleDegrees[8, 54, "totalLocal"]
CoordinateLocalCondensateComparisonToMomentum[8, 54]
CoordinateLocalCondensateComparisonToMomentum[8, 54,
  $BcCoordinateDefaultParameters, "IncludeG3" -> True]
\end{verbatim}
The labels mean
\begin{align}
  G2c &: S_c^{(G^2)}(-x)S_b^{(0)}(x),&
  G2b &: S_c^{(0)}(-x)S_b^{(G^2)}(x),\\
  G3c &: S_c^{(G^3)}(-x)S_b^{(0)}(x),&
  G3b &: S_c^{(0)}(-x)S_b^{(G^3)}(x).
\end{align}
At \(M^2=8~\mathrm{GeV}^2\), \(s_0=54~\mathrm{GeV}^2\), the local-coordinate
test gives
\begin{equation}
  \theta_{\rm pert+G2local}=43.321950845^\circ,\qquad
  \theta_{\rm pert+G2local+G3local}=43.308150036^\circ.
\end{equation}
The word ``local'' is important.  It means that the condensate is inserted on
one heavy-quark line at a time.  For \(G^3\), this reproduces the
single-line momentum-space \(G^3\) implementation used here.  For \(G^2\), it
does not yet reproduce the complete dimension--4 result, because there is also
an open-field term in which one background gluon is attached to the charm line
and one to the bottom line.

Numerically, for the \(AA\) channel at the same test point,
\begin{align}
  \Pi_{G^2,\mathrm{local}}^{AA}
  &=
  -3.1655383105\times10^{-4},\\
  \Pi_{G^2,\mathrm{full\ momentum}}^{AA}
  &=
  -3.6293326185\times10^{-4},\\
  \Pi_{G^2,\mathrm{residual}}^{AA}
  &=
  -4.6379430799\times10^{-5}.
\end{align}
This residual is the target for the missing coordinate-space
\(S_c^{(G)}(-x)S_b^{(G)}(x)\) derivation.  By contrast, the local coordinate
and momentum-space \(G^3\) moments agree numerically at this point:
\begin{align}
  \Pi_{G^3,\mathrm{local}}^{AA}
  &=
  6.8538572092\times10^{-5},\\
  \Pi_{G^3,\mathrm{momentum}}^{AA}
  &=
  6.8538572092\times10^{-5}.
\end{align}
This should not yet be quoted as the complete condensate-corrected coordinate
result.  The reason is that the full dimension--4 contribution also contains
the cross-line \(S_c^{(G)}(-x)S_b^{(G)}(x)\) term.  The helper
\texttt{CoordinateLocalCondensateComparisonToMomentum} reports the difference
between the full momentum-space \(G^2\) moment and the local coordinate
\(G2c+G2b\) moment; this residual is the piece that must be reproduced from
the open-field coordinate-space propagator.  The \(G^3\) comparison is
optional because it is slower; setting \texttt{"IncludeG3" -> True} checks
that the local coordinate \(G3c+G3b\) integration matches the momentum-space
\(G^3\) contribution.

\section{Checklist for independent verification}

A theorist can check the calculation in the following order:
\begin{enumerate}
  \item Confirm the current convention, especially the explicit factor of
  $i$ in $J^B_\mu$, the normalization $p^\alpha/(m_b+m_c)$, and the
  definition of $\sigma_{\mu\nu}$.
  \item Reproduce the three projected perturbative traces and verify the
  numerator functions $N^{AA}(s)$, $N^{AB}(s)$, and $N^{BB}(s)$.
  \item Check the two-body phase-space normalization
  $N_c\sqrt{\lambda}/(16\pi^2s)$ used in $\rho_{\rm pert}^{ij}(s)$.
  \item Verify the background-field propagator terms and the color factor in
  the $S^{(G)}S^{(G)}$ contribution.
  \item Independently derive the direct Borel formula for the $G^2$ and $G^3$
  terms, including the $-i$ phase convention.
  \item Check whether the standard single-line $G^3$ propagator contribution
  is sufficient, or whether cross-line open-field $G^3$ terms must be added.
  \item Check that $\Pi_{G^2}/\Pi_{\rm pert}$ and
  $\Pi_{G^3}/\Pi_{\rm pert}$ remain small in the accepted Borel window.
  \item Use the $K_{ij}$ sensitivity model only as an estimate of possible
  channel-dependent perturbative $O(\alpha_s)$ uncertainty.
  \item Derive the true $O(\alpha_s)$ two-loop perturbative spectral
  densities separately.
  \item Determine the final working window using OPE convergence, Borel
  stability, and pole-dominance criteria.
\end{enumerate}

\end{document}
